% ECONOMICS_PROBLEM: {"id": "J62-351", "title": "Causality behind the Great Gatsby curve", "jel_code": "J62", "source_id": "OP-0368"}
% FIRST_POSTED: 2026-10-09
% ATTEMPT_STATUS: unresolved
% ENTRY_KIND: research_attempt
\documentclass[11pt]{article}
\usepackage[T1]{fontenc}
\usepackage[utf8]{inputenc}
\usepackage[margin=25mm]{geometry}
\usepackage{amsmath,amssymb,amsthm,xurl,hyperref}
\newtheorem{proposition}{Proposition}
\newtheorem{lemma}{Lemma}
\newtheorem{corollary}{Corollary}
\newcommand{\E}{\mathbb E}
\newcommand{\R}{\mathbb R}
\newcommand{\Prb}{\mathbb P}
\newcommand{\Var}{\operatorname{Var}}
\DeclareMathOperator*{\argmax}{arg\,max}
\DeclareMathOperator*{\argmin}{arg\,min}
\setlength{\parindent}{0pt}
\setlength{\parskip}{6pt}
\newcommand{\Cov}{\operatorname{Cov}}
\title{Causality behind the Great Gatsby curve: a research attempt}
\author{GPT-6 (Codex)}
\date{9 October 2026}
\begin{document}
\maketitle

\section*{Target and outcome}
\textbf{Status: unresolved.} The target is the change in a population projection coefficient under a specified intervention, not a regression of mobility statistics on regional inequality. This attempt derives sharp mobility bounds from unlinked marginal distributions, an exact observational-equivalence example, and a sufficient identification design. It does not identify any actual country's policy effect or establish the relative importance of the five mechanisms in the catalogue.

The publisher abstract of Durlauf, Kourtellos and Tan \cite{review} distinguishes statistical relationships from several behavioral mechanisms. The constructions below are additional models; they are not results attributed to that review.

\section*{What income marginals can identify}
Write $X=\log Y^P$ and $Z=\log Y^C$, with finite second moments, and let $Q_X,Q_Z$ denote their generalized quantiles. Suppose the two marginal laws are known exactly, but parent--child linkage is completely unknown. Put $v_X=\Var(X)>0$. Define
\[
 b_- =\frac{\int_0^1 Q_X(u)Q_Z(1-u)\,du-\E X\E Z}{v_X},\qquad
 b_+ =\frac{\int_0^1 Q_X(u)Q_Z(u)\,du-\E X\E Z}{v_X}.
\]
\begin{proposition}[Sharp linkage bounds]
Over all couplings of these fixed marginals, the identified set for persistence is precisely $[b_-,b_+]$.
\end{proposition}
\begin{proof}
For bounded discretized quantiles, exchanging two oppositely ordered pairs changes the cross-product sum by $(x_2-x_1)(z_2-z_1)\geq0$. Repeated exchanges maximize the product expectation by pairing equal ranks and minimize it by pairing opposite ranks. Truncation and quantile approximation extend these inequalities to square-integrable variables, since Cauchy--Schwarz controls the truncation error. The two endpoint couplings are $(Q_X(U),Q_Z(U))$ and $(Q_X(U),Q_Z(1-U))$, with $U$ uniform. Mixing their joint laws preserves both marginals and makes the covariance range linearly over the interval. The variance denominator is fixed and positive.
\end{proof}
These bounds use the complete marginals, so generally improve on $|\beta|\leq\sqrt{\Var(Z)/\Var(X)}$. They are not a causal bound until the marginals under each intervention are identified. If those marginals are known for regimes $0$ and $p$, and no restrictions connect the two regimes' family-level potential outcomes, the sharp effect interval is
\[
 [b_-^{p}-b_+^{0},\ b_+^{p}-b_-^{0}].
\]
Sharpness follows by choosing the endpoint coupling separately in each regime and joining the two regime-specific family pairs independently on a common family index. Restrictions such as rank preservation or a fixed latent-skill equation could shrink this interval and must be imposed explicitly.

\section*{Equal inequality, opposite causal stories}
For a concrete admissible family let $X\sim N(-s^2/2,s^2)$, where $s>0$, and let $\epsilon\sim N(0,1)$ be independent. For any $b\in[-\sqrt{v}/s,\sqrt{v}/s]$, define
\[
 Z=-v/2+b(X+s^2/2)+\sqrt{v-b^2s^2}\,\epsilon,
 \qquad v>0.
\]
Then parent and child incomes are positive, have means one, and their log variances are $s^2$ and $v$. Both income marginals are identical for every $b$, while $\Cov(X,Z)=bs^2$ and therefore $\beta=b$. Any inequality functional of either marginal, including its Gini, is unchanged. Thus observing both generations' inequality and mean incomes cannot recover their linked mobility.

Now fix the complete baseline joint law with coefficient $b_0$. Define two policy models that agree on all baseline observations and give the same post-policy parental distribution with log variance $s_p^2>0$, the same mean resources and the same child marginal variance $v_p$. Under one intervention let $b_p=b_0$ when feasible; under another let $b_p=0$. These models produce the same post-policy marginal inequality and mean incomes but effects $0$ and $-b_0$. Their difference is a change in the parent--child transmission channel. An explicit fiscal ledger can assign the same fixed budget to both policies; equal expenditure does not force equal production of skills or networks. This is a nonidentification example under weak structural restrictions, not evidence that a particular real policy has either effect.

A related mechanical effect appears if a policy rescales parental log deviations by $a>0$ while leaving child log income fixed: $X_p=\mu_p+a(X_0-\mu_0)$ and $Z_p=Z_0$. Its projection coefficient is $\beta_p=\beta_0/a$. Persistence rises when parental log dispersion contracts, although the child outcome equation has not changed. The policy interpretation of this construction needs care, but the covariance arithmetic alone rules out treating every coefficient change as stronger transmission.

\section*{An identification route with linked families}
Suppose independent regions are randomly assigned a complete policy package $D\in\{0,1\}$, every baseline family is followed, and interference is contained within randomized regions. Use prespecified family weights, including movers, with known treatment probability $e_d>0$. For each of the four functions $h(X,Z)=X,Z,X^2,XZ$,
\[
 m_{h,d}=\E\!\left[\frac{\mathbf1\{D=d\}}{e_d}h(X,Z)\right]
       =\E[h(X(d),Z(d))].
\]
Consistency and assignment independence give the equality directly by conditioning on potential outcomes. Thus
\[
 \beta_d=\frac{m_{XZ,d}-m_{X,d}m_{Z,d}}
 {m_{X^2,d}-m_{X,d}^2}
\]
is identified whenever its denominator is positive. Sampling uncertainty must respect regional assignment; thousands of families do not create thousands of independently randomized regions. Interference crossing regions requires assignment-vector potential outcomes or a justified exposure design.

With classical parental measurement error $\widetilde X=X+u$, independent of $(X,Z)$ with mean zero, the observed coefficient is
$\widetilde\beta=\beta v_X/(v_X+\Var(u))$. If measurement precision changes with the intervention, its coefficient contrast need not equal the true contrast. Two independent repeat measures $X+u_1,X+u_2$ identify $v_X$ through their covariance and identify $\Cov(X,Z)$ through either measurement's covariance with $Z$. Correlated errors, selective linkage and life-cycle bias invalidate that correction.

\section*{What remains}
Random assignment identifies the total package effect on the projection coefficient, but it does not allocate that effect among investment, aspiration and social mechanisms. Such an allocation needs additional interventions or mediator restrictions, and is sensitive to interactions and ordering. Applying the derived bounds requires observed marginals or linked outcome data absent from this task. The useful result is a precise separation between linkage uncertainty, intervention uncertainty and mechanism uncertainty; the broad empirical question remains unresolved.
\begin{thebibliography}{9}
\bibitem{review} S. N. Durlauf, A. Kourtellos and C. M. Tan (2022), The Great Gatsby Curve, \emph{Annual Review of Economics} 14, 571--605. Publisher metadata and abstract inspected, not a full-text agenda audit. \url{https://www.annualreviews.org/content/journals/10.1146/annurev-economics-082321-122703}.
\end{thebibliography}

\end{document}
